% SEGMENT OLP-0534-S01
% Part: set-theory
% Chapter: story
% Section: cumulative-approach

\documentclass[../../../include/open-logic-section]{subfiles}

\begin{document}

\olfileid{sth}{story}{approach}
\olsection{திரள்வுறும் படிநிலைமுறை அணுகுமுறை}

கணங்களைப் படிநிலைகளாக நாம் எவ்வாறு வகைப்படுத்துவோம் என்பதன் சற்றே
விரிவான கூற்று இதோ:
\begin{quote}
	கணங்கள் \emph{படிநிலைகளில்} உருவாகின்றன. ஒவ்வொரு படிநிலை~$S$ க்கும்,
	$S$ க்கு \emph{முந்திய} சில படிநிலைகள் உள்ளன. படிநிலை~$S$ இல், $S$
	க்கு முந்திய படிநிலைகளில் உருவான கணங்களைக் கொண்ட ஒவ்வொரு தொகுப்பும்
	ஒரு கணமாக உருவாக்கப்படுகிறது. படிநிலைகளில் உருவாக்கப்படும் கணங்களைத்
	தவிர வேறு கணங்கள் இல்லை. \citep[p.~323]{Shoenfield:AST}
\end{quote}
இது \emph{திரள்வுறும் படிநிலைமுறைக் கணக் கருத்தாக்கத்தின்} ஒரு
வரைவுரு. \olref[sth][][]{part} இல் நாம் வழங்கும் முறைசார் கணக்
கோட்பாட்டிற்கு இதுவே அடித்தளமாகும்.

இதை இன்னும் சற்று விரிவாக ஆராய்வோம். ஷோன்ஃபீல்ட் இந்தச் செயல்முறையை
விவரிப்பதுபோல், ஒவ்வொரு படிநிலையிலும், முந்திய படிநிலைகளிலிருந்து நமக்குக்
கிடைத்த {கணங்களை} கொண்டு புதிய கணங்களை உருவாக்குகிறோம். எனவே
ஷோன்ஃபீல்டின் சித்திரத்தில், தொடக்கப் படிநிலையான படிநிலை~$0$ க்கு
\emph{முந்திய} படிநிலைகள் இல்லை; ஆகவே \emph{மேலும் உறுதியாக}, முந்திய
படிநிலைகளிலிருந்து நமக்குக் கிடைக்கும் கணங்களும் இல்லை.\footnote{ஒரு
முதல் படிநிலை \emph{உள்ளது} என்று நாம் ஏன் எடுகோள வேண்டும்?
\olref[z][story]{sec} இலுள்ள \stagesord{} இன் அடிக்குறிப்பைப் பார்க்கவும்.}
எனவே ஒரே ஒரு கணத்தை உருவாக்குகிறோம்: !!{element}s எதுவுமில்லாத
கணம்~$\emptyset$. படிநிலை~$1$ இல், முந்திய படிநிலைகளிலிருந்து துல்லியமாக
ஒரே ஒரு கணம் நமக்குக் கிடைக்கிறது; ஆகவே ஒரே புதிய கணம்
$\{\emptyset\}$ உருவாகிறது. படிநிலை~$2$ இல், முந்திய படிநிலைகளிலிருந்து
இரண்டு கணங்கள் நமக்குக் கிடைக்கின்றன; $\{\{\emptyset\}\}$ மற்றும்
$\{\emptyset, \{\emptyset\}\}$ ஆகிய இரு புதிய கணங்களை உருவாக்குகிறோம்.
படிநிலை~$3$ இல், முந்திய படிநிலைகளிலிருந்து நான்கு கணங்கள் நமக்குக்
கிடைக்கின்றன; எனவே பன்னிரண்டு புதிய கணங்களை உருவாக்குகிறோம்\ldots.
ஆகவே கணங்களின் திரள்வுறும் படிநிலைமுறைச் சித்திரம் இதைப் போன்றிருக்கும்
(எண்கள் படிநிலைகளைக் குறிக்கின்றன):
\begin{center}
	\begin{tikzpicture}[scale=0.6]
	\tikzset{cut_here/.style={densely dotted}}
	\draw (-4,6) -- (-1, 0)-- (2,6); 
	\draw[cut_here] (-5,8)--(-4,6);
	\draw[cut_here] (2,6)--(3,8);
	\draw(-1.5, 1)--(-0.5, 1);
	\draw(-2, 2)--(0, 2);
	\draw(-2.5, 3)--(0.5,3);
	\draw(-3, 4)--(1, 4);
	\draw(-3.5, 5)--(1.5, 5);
	\draw(-4, 6)--(2, 6);
	\node[label] at (0, 0) {\small 0};
	\node[label] at (0.5, 1) {\small 1};
	\node[label] at (1, 2) {\small 2};
	\node[label] at (1.5, 3) {\small 3};
	\node[label] at (2, 4) {\small 4};
	\node[label] at (2.5, 5) {\small 5};
	\node[label] at (3, 6) {\small 6};
	\end{tikzpicture}
\end{center}
அப்படியானால், இந்தக் கதையை நாம் ஏன் ஏற்க வேண்டும்?

ஒரு காரணம், இது நேர்த்தியானதும் கையாளத்தக்கதுமான கதை. மிக வெளிப்படையான
கதையாகிய முறைசாராப் பண்புவழிக் கணமாக்கல் வீழ்ந்துவிட்டதால், நேர்த்தியான
ஏதோவொன்று நமக்குத் தேவைப்படுகிறது.

ஆனால் இந்தக் கதை \emph{வெறுமனே} நேர்த்தியானது மட்டுமல்ல. இந்தக் கதையை
அடிப்படையாகக் கொண்ட எந்தக் கணக் கோட்பாடும் \emph{முரணற்றதாக} இருக்கும்
என்று நம்ப நமக்கு நல்ல காரணம் உண்டு. அதன் காரணம் இதோ.

திரள்வுறும் படிநிலைமுறைக் கணக் கருத்தாக்கத்தின்படி, கணங்களைப்
படிநிலைகளில் உருவாக்குகிறோம்; அவற்றின் !!{element}s \emph{ஏற்கெனவே}
கிடைத்த பொருள்களாக இருக்க வேண்டும். ஆகவே, எந்தப் படிநிலை~$S$ க்கும்,
பின்வரும் கணத்தை உருவாக்கலாம்:
\[
	R_S = \Setabs{x}{x \notin x \text{ மற்றும் $x$, $S$ க்கு முன்பு கிடைத்தது}}
\]
ரஸ்ஸலின் முரண்பாட்டை நிறுவுவதில் இடம்பெறும் பகுத்தறிதல், $R_S$ என்பதே
படிநிலை~$S$ க்கு முன்பு கிடைக்கவில்லை என்பதை இப்போது நிறுவும். அது ஒரு
முரண் அல்ல. மேலும், திரள்வுறும் படிநிலைமுறைக் கணக் கருத்தாக்கத்தை நாம்
ஏற்றுக்கொண்டால், ரஸ்ஸல் கணத்தையே உருவாக்க முடியும் என்று நாம்
\emph{எதிர்பார்த்திருக்கவும்} கூடாது. ஏனெனில் அது “எப்போதாவது கிடைக்கப்
போகும்” தன்னுறுப்பில்லாத எல்லாக் கணங்களின் கணமாக இருக்கும். சுருக்கமாக:
ரஸ்ஸல் கணத்தை நம்மால் (மெய்ப்பிக்கத்தக்கவாறு) உருவாக்க முடியவில்லை என்பது,
திரள்வுறும் படிநிலைமுறைக் கதையைக் கொள்கையில் \emph{வியப்புக்குரியதல்ல};
அதைத்தான் நாம் \emph{முன்கணிப்போம்}.

%ஒரு பொருளில், திரள்வுறும் படிநிலைமுறைக் கணக் கருத்தாக்கம் ரஸ்ஸலின் முரண்பாட்டிற்கு \emph{முன்கணிப்புவாதியின்} பதிலைப் போன்ற ஒரு பதிலைத் தருகிறது. தன்னுறுப்பில்லாத கணங்கள்$_n$ அனைத்தின் தொகுப்பு ஒரு கணம்$_{n+1}$ என்றும், இந்தக் கணம்$_n$ தன்னுறுப்புடையதாக இருக்காது என்றும் முன்கணிப்புவாதி கூறினார். (உண்மையில், ஒரு கணம் தன்னுறுப்புடையதா என்று கேட்க முயல்வதையே பெரும்பாலான முன்கணிப்புவாதிகள் \emph{இலக்கணமற்றதாகக்} கருதுவர்.)
%
%ஆனால் திரள்வுறும் படிநிலைமுறை அணுகுமுறைக்கும் முன்கணிப்புவாத அணுகுமுறைக்கும் ஒரு முக்கிய வேறுபாடு உள்ளது: திரள்வுறும் படிநிலைமுறை அணுகுமுறை நமது பொருள்கள் அனைத்தையும் \emph{ஒரே வகையைச் சேர்ந்தவையாகக்} கருதுகிறது. முன்கணிப்புவாதி வெற்றுக் கணம்$_0$ ஒன்றையும் (அதாவது, !!{element}s இல்லாத ஒரு கணம்$_0$), வெற்றுக் கணம்$_1$ ஒன்றையும் (அதாவது, !!{element}s இல்லாத ஒரு கணம்$_1$) கொண்டிருப்பார்; இவை \emph{வேறுபட்ட பொருள்களாக} இருக்கும். ஆனால் திரள்வுறும் படிநிலைமுறை அணுகுமுறையில், $\emptyset$ என்ற ஒரே ஒரு வெற்றுக் கணம் மட்டுமே உள்ளது; அது படிநிலை வரிசையின் ஒவ்வொரு படிநிலையிலும் கிடைக்கும்.

%உண்மையில், இந்த அத்தியாயத்தின் தொடக்கத்தில் கூறப்பட்ட உறுப்புசார் சமத்துவ அடிகோள், இந்தப் படிநிலை வரிசையிலுள்ள உறுப்புகளுக்கு மெய்யாக இருக்கும் (இதனால் ஒரே ஒரு “வெற்றுக்” கணம் \emph{மட்டுமே} இருக்க முடியும்). ஆனால் உறுப்புசார் சமத்துவத்தை நியாயப்படுத்த திரள்வுறும் படிநிலைமுறைக் கருத்தாக்கத்தைப் பயன்படுத்த நான் கருதவில்லை (\cite{Boolos1971} ஐப் பார்க்கவும்). மீண்டும்: உறுப்புகளால் நிர்ணயிக்கப்படும் தொகுப்புகளிலேயே நமக்கு ஆர்வம் இருப்பதால் உறுப்புசார் சமத்துவத்தை ஏற்க வேண்டும் எனக் கூறுகிறேன்.

\end{document}
